\documentclass[10pt,a4paper]{amsart}
\usepackage{amsmath,amssymb,amsthm,graphicx}
\usepackage{verbatim}
\newenvironment{reference}{\comment}{\endcomment}
%\newenvironment{reference}{}{}
\newenvironment{slogan}{\comment}{\endcomment}
\newenvironment{history}{\comment}{\endcomment}

% For commutative diagrams we use Xy-pic
\usepackage[all]{xy}

% We use 2cell for 2-commutative diagrams.
\xyoption{2cell}
% Theorem environments.
%
\theoremstyle{plain}
\newtheorem{theorem}[subsection]{Theorem}
\newtheorem{proposition}[subsection]{Proposition}
\newtheorem{lemma}[subsection]{Lemma}

\theoremstyle{definition}
\newtheorem{definition}[subsection]{Definition}
\newtheorem{example}[subsection]{Example}
\newtheorem{exercise}[subsection]{Exercise}
\newtheorem{situation}[subsection]{Situation}

\theoremstyle{remark}
\newtheorem{remark}[subsection]{Remark}
\newtheorem{remarks}[subsection]{Remarks}

\numberwithin{equation}{subsection}

% Macros
%
\def\lim{\mathop{\mathrm{lim}}\nolimits}
\def\colim{\mathop{\mathrm{colim}}\nolimits}
\def\Spec{\mathop{\mathrm{Spec}}}
\def\Hom{\mathop{\mathrm{Hom}}\nolimits}
\def\Ext{\mathop{\mathrm{Ext}}\nolimits}
\def\SheafHom{\mathop{\mathcal{H}\!\mathit{om}}\nolimits}
\def\SheafExt{\mathop{\mathcal{E}\!\mathit{xt}}\nolimits}
\def\Sch{\mathit{Sch}}
\def\Mor{\mathop{\mathrm{Mor}}\nolimits}
\def\Ob{\mathop{\mathrm{Ob}}\nolimits}
\def\Sh{\mathop{\mathit{Sh}}\nolimits}
\def\NL{\mathop{N\!L}\nolimits}
\def\CH{\mathop{\mathrm{CH}}\nolimits}
\def\proetale{{pro\text{-}\acute{e}tale}}
\def\etale{{\acute{e}tale}}
\def\QCoh{\mathit{QCoh}}
\def\Ker{\mathop{\mathrm{Ker}}}
\def\Im{\mathop{\mathrm{Im}}}
\def\Coker{\mathop{\mathrm{Coker}}}
\def\Coim{\mathop{\mathrm{Coim}}}

% Boxtimes
%
\DeclareMathSymbol{\boxtimes}{\mathbin}{AMSa}{"02}

%
% Macros for moduli stacks/spaces
%
\def\QCohstack{\mathcal{QC}\!\mathit{oh}}
\def\Cohstack{\mathcal{C}\!\mathit{oh}}
\def\Spacesstack{\mathcal{S}\!\mathit{paces}}
\def\Quotfunctor{\mathrm{Quot}}
\def\Hilbfunctor{\mathrm{Hilb}}
\def\Curvesstack{\mathcal{C}\!\mathit{urves}}
\def\Polarizedstack{\mathcal{P}\!\mathit{olarized}}
\def\Complexesstack{\mathcal{C}\!\mathit{omplexes}}
% \Pic is the operator that assigns to X its picard group, usage \Pic(X)
% \Picardstack_{X/B} denotes the Picard stack of X over B
% \Picardfunctor_{X/B} denotes the Picard functor of X over B
\def\Pic{\mathop{\mathrm{Pic}}\nolimits}
\def\Picardstack{\mathcal{P}\!\mathit{ic}}
\def\Picardfunctor{\mathrm{Pic}}
\def\Deformationcategory{\mathcal{D}\!\mathit{ef}}
\usepackage{cleveref}
\crefname{originalreference}{original reference}{original references}
\Crefname{originalreference}{Original reference}{Original references}
\makeatletter
\newlabel{spaces-topologies-lemma-fppf-covering-surjective}{{0469}{source}{}{sourcepointer-1}{}}
\newlabel{spaces-topologies-lemma-fppf-covering-surjective@cref}{{[originalreference][0][]0469}{[1][0][]source}{}{}{}}
\makeatother
\crefname{originallemma}{original lemma}{original lemmas}
\Crefname{originallemma}{Original lemma}{Original lemmas}
\makeatletter
\newlabel{lemma-curiosity}{{06NR}{source}{}{sourcepointer-2}{}}
\newlabel{lemma-curiosity@cref}{{[originallemma][0][]06NR}{[1][0][]source}{}{}{}}
\makeatother
\crefname{originalreference}{original reference}{original references}
\Crefname{originalreference}{Original reference}{Original references}
\makeatletter
\newlabel{spaces-morphisms-lemma-finite-type-local}{{040Y}{source}{}{sourcepointer-3}{}}
\newlabel{spaces-morphisms-lemma-finite-type-local@cref}{{[originalreference][0][]040Y}{[1][0][]source}{}{}{}}
\makeatother
\crefname{originalreference}{original reference}{original references}
\Crefname{originalreference}{Original reference}{Original references}
\makeatletter
\newlabel{spaces-morphisms-lemma-finite-presentation-local}{{0410}{source}{}{sourcepointer-4}{}}
\newlabel{spaces-morphisms-lemma-finite-presentation-local@cref}{{[originalreference][0][]0410}{[1][0][]source}{}{}{}}
\makeatother
\crefname{originalreference}{original reference}{original references}
\Crefname{originalreference}{Original reference}{Original references}
\makeatletter
\newlabel{spaces-morphisms-lemma-fppf-open}{{042S}{source}{}{sourcepointer-5}{}}
\newlabel{spaces-morphisms-lemma-fppf-open@cref}{{[originalreference][0][]042S}{[1][0][]source}{}{}{}}
\makeatother
\crefname{originalreference}{original reference}{original references}
\Crefname{originalreference}{Original reference}{Original references}
\makeatletter
\newlabel{descent-lemma-flat-finitely-presented-permanence}{{02KL}{source}{}{sourcepointer-6}{}}
\newlabel{descent-lemma-flat-finitely-presented-permanence@cref}{{[originalreference][0][]02KL}{[1][0][]source}{}{}{}}
\makeatother
\title{Permanence of local finite presentation/type along fppf covers with surjective flat lfp \$f\$}
\author{The Stacks Project Authors}
\date{stacks-project:spaces-descent.tex:local-snapshot-20261012; GNU FDL 1.2 or later}
\begin{document}
\maketitle
\noindent\textit{Editorial note.} All delivered mathematical text is the author's own, copied verbatim from the pinned Stacks Project chapter file spaces-descent.tex; no statement or proof step was rewritten, reordered or removed.\par\medskip
\noindent\textit{Editorial note.} The counted span cites three Morphisms of Spaces lemmas and one Descent lemma, recorded in assembly.external\_references with canonical tags as navigation pointers.\par\medskip
\noindent\textit{Allowed standard prerequisites (editorial).}\begin{itemize}
\item Local finite type/presentation is \textbackslash{}'etale local on the source. Stacks Project, Morphisms of Spaces, Lemmas spaces-morphisms-lemma-finite-type-local and spaces-morphisms-lemma-finite-presentation-local. Justifies the reduction to affine \$B\$ and \$Y\$.
\item Openness of fppf morphisms on points. Stacks Project, Morphisms of Spaces, Lemma spaces-morphisms-lemma-fppf-open. Lets the proof pick \$U\$ with \$U \textbackslash{}to Y\$ surjective.
\item Scheme-side permanence. Stacks Project, Descent, Lemma descent-lemma-flat-finitely-presented-permanence. Concludes for the scheme diagram.
\end{itemize}\medskip
\ifdefined\hypertarget\hypertarget{sourcepointer-1}{}\fi
\noindent\textit{Source reference.} Original reference 0469: sources/stacks-chow/tags.txt, source lines 5464--5464. Cross-chapter standard prerequisite cited by the author inside the counted span; it is used as a pointer and its own proof is not part of this item.\par\medskip
\ifdefined\hypertarget\hypertarget{sourcepointer-2}{}\fi
\noindent\textit{Source reference.} Original lemma 06NR: sources/stacks-chow/tags.txt, source lines 8545--8545. Cross-chapter standard prerequisite cited by the author inside the counted span; it is used as a pointer and its own proof is not part of this item.\par\medskip
\ifdefined\hypertarget\hypertarget{sourcepointer-3}{}\fi
\noindent\textit{Source reference.} Original reference 040Y: sources/stacks-chow/tags.txt, source lines 5273--5273. Cross-chapter standard prerequisite cited by the author inside the counted span; it is used as a pointer and its own proof is not part of this item.\par\medskip
\ifdefined\hypertarget\hypertarget{sourcepointer-4}{}\fi
\noindent\textit{Source reference.} Original reference 0410: sources/stacks-chow/tags.txt, source lines 5275--5275. Cross-chapter standard prerequisite cited by the author inside the counted span; it is used as a pointer and its own proof is not part of this item.\par\medskip
\ifdefined\hypertarget\hypertarget{sourcepointer-5}{}\fi
\noindent\textit{Source reference.} Original reference 042S: sources/stacks-chow/tags.txt, source lines 5339--5339. Cross-chapter standard prerequisite cited by the author inside the counted span; it is used as a pointer and its own proof is not part of this item.\par\medskip
\ifdefined\hypertarget\hypertarget{sourcepointer-6}{}\fi
\noindent\textit{Source reference.} Original reference 02KL: sources/stacks-chow/tags.txt, source lines 3376--3376. Cross-chapter standard prerequisite cited by the author inside the counted span; it is used as a pointer and its own proof is not part of this item.\par\medskip
\begin{lemma}
\label{lemma-flat-finitely-presented-permanence}
Let $S$ be a scheme. Let
$$
\xymatrix{
X \ar[rr]_f \ar[rd]_p & &
Y \ar[dl]^q \\
& B
}
$$
be a commutative diagram of morphisms of algebraic spaces over $S$.
Assume that $f$ is surjective, flat, and locally of finite presentation
and assume that $p$ is locally of finite presentation (resp.\ locally
of finite type). Then $q$ is locally of finite presentation
(resp.\ locally of finite type).
\end{lemma}

\begin{proof}
Since $\{X \to Y\}$ is an fppf covering, it induces a surjection of
fppf sheaves (Topologies on Spaces, Lemma
\ref{spaces-topologies-lemma-fppf-covering-surjective}) and the
lemma is a special case of Lemma \ref{lemma-curiosity}.
On the other hand, an easier argument is to deduce it from
the analogue for schemes. Namely, the problem is \'etale local
on $B$ and $Y$ (Morphisms of Spaces, Lemmas
\ref{spaces-morphisms-lemma-finite-type-local} and
\ref{spaces-morphisms-lemma-finite-presentation-local}).
Hence we may assume that $B$ and $Y$ are affine
schemes. Since $|X| \to |Y|$ is open
(Morphisms of Spaces, Lemma \ref{spaces-morphisms-lemma-fppf-open}),
we can choose an affine
scheme $U$ and an \'etale morphism $U \to X$ such that the
composition $U \to Y$ is surjective. In this case the result
follows from Descent, Lemma
\ref{descent-lemma-flat-finitely-presented-permanence}.
\end{proof}

\end{document}
